% !TEX program = xelatex
% 问题一分析——思维导图（按原图复刻，XeLaTeX 编译）
\documentclass[border=6pt]{standalone} % border=6pt：PDF 页面四周留白宽度
\usepackage{ctex}
\usepackage{fontspec}
\setmainfont{Times New Roman}
\usepackage{amsmath,amssymb}
\usepackage{tikz}
\usetikzlibrary{arrows.meta,decorations.pathreplacing}

% ---------- 取自原图的配色 ----------
\definecolor{dgreen}{HTML}{0E4F43}  % 深绿边框/虚线
\definecolor{pOne}{HTML}{EAF6EC}    % Step1 面板底色（浅绿）
\definecolor{pTwo}{HTML}{E7F0FA}    % Step2 面板底色（浅蓝）
\definecolor{pThr}{HTML}{F6E9D4}    % Step3 面板底色（浅橙）
\definecolor{pFou}{HTML}{DEF2F7}    % Step4 面板底色（浅青）
\definecolor{pinkA}{HTML}{F8F1F7}   % 免检/复核区域（浅粉白）
\definecolor{pinkB}{HTML}{F7D8DC}   % 窗口统计量（粉红）
\definecolor{cyanA}{HTML}{8CDDE6}   % 首轮/复核异常（青）
\definecolor{whiteA}{HTML}{F8F5FA}  % 公式框（浅紫白）
\definecolor{cyanB}{HTML}{8FE0EA}   % Step2 子面板（青）
\definecolor{cyanC}{HTML}{DFF6F8}   % Step3 增益/损伤/季节（浅青）
\definecolor{greenC}{HTML}{EBF7ED}  % Step3 中维护/大维护（浅绿）

\begin{document}
	\begin{tikzpicture}[font=\large, % 全图节点文字字号（改 \normalsize 可整体缩小文字）
		% --- 箭头样式：line width=1.6pt 为箭头线粗；length=6pt 为箭头头部长度；width=5pt 为箭头头部宽度 ---
		arr/.style={line width=1.6pt,-{Stealth[length=6pt,width=5pt]}},
		% --- 无箭头连接线：line width=1.6pt 为线粗 ---
		pln/.style={line width=1.6pt},
		% --- 圆角矩形节点：rounded corners=5pt 圆角半径；line width=2pt 边框宽度；
		%     inner sep=2.5mm 文字到边框的内边距；minimum height=11mm 盒子最小高度（宽度随文字自动伸展）---
		boxR/.style={rounded corners=5pt,draw=dgreen,line width=2pt,
			inner sep=2.5mm,minimum height=11mm,align=center},
		% --- 直角矩形节点：line width=2pt 边框宽度；inner sep、minimum height 含义同上 ---
		boxS/.style={draw=dgreen,line width=2pt,
			inner sep=2.5mm,minimum height=11mm,align=center},
		% --- 主面板虚线边框：line width=1pt 虚线线粗 ---
		panel/.style={draw=dgreen,dotted,line width=2pt},
		% --- Step2 青色子面板：rounded corners=5pt 圆角半径；line width=2pt 边框宽度 ---
		subp/.style={rounded corners=5pt,draw=dgreen,line width=2pt},
		]
		
		% ================= 四个虚线面板 =================
		% rectangle 两坐标为左下角与右上角，单位 cm
		\draw[panel,fill=pOne] (0,8.1)    rectangle (23,16.3);   % Step1 面板：宽 23.0，高 8.2
		\draw[panel,fill=pTwo] (23.2,0.1) rectangle (34,16.3);   % Step2 面板：宽 10.8，高 16.2
		\draw[panel,fill=pThr] (12.1,0.1) rectangle (22.9,7.6);  % Step3 面板：宽 10.8，高 7.5
		\draw[panel,fill=pFou] (0,0.1)    rectangle (11.6,7.6);  % Step4 面板：宽 11.6，高 7.5
		
		% 面板标题（anchor=north west：所给坐标为文本块左上角）
		\node[anchor=north west,align=left] at (0.7,16.0)  {\textbf{Step 1：数据预处理}\\(Hample滤波器)};
		\node[anchor=north west,align=left] at (23.8,16.0) {\textbf{Step 2：分析数据特征}\\(STL时间序列分解\quad FFT\quad ACF)};
		\node[anchor=north west] at (12.6,7.35) {\textbf{Step 3：维护的影响}};
		\node[anchor=north west] at (0.6,7.35)  {\textbf{Step 4：透水率指标}};
		
		% ================= Step 1：数据预处理 =================
		% 各 at (x,y) 为盒子中心坐标
		\node[boxR,fill=pinkA] (A1) at (7.7,15.4)  {免检区域 $\mathrm{flag}_i^{(1)}$};
		\node[boxR,fill=pinkA] (A2) at (7.7,14.0)  {复核区域 $\mathrm{flag}_i^{(2)}$};
		\node[boxR,fill=cyanA] (B1) at (13.9,15.4) {首轮异常 $\mathcal{O}_1$};
		\node[boxR,fill=cyanA] (B2) at (13.9,14.0) {复核异常 $\mathcal{O}_2$};
		\node[boxS,fill=whiteA] (E) at (19.1,14.7) {异常点集\\$\mathcal{O}=\mathcal{O}_1\cup\mathcal{O}_2$}; % 双行盒子，高度自动叠加
		
		\node[boxR,fill=pinkB] (W1) at (8.1,12.4)  {动态时间窗口\quad $\mathcal{N}(t_i)$};
		\node[boxR,fill=pinkB] (W2) at (8.1,10.8)  {滑动窗口的中位数 $m_i$};
		\node[boxR,fill=pinkB] (W3) at (8.1,9.2)   {中值绝对偏差 $\mathrm{MAD}_i$};
		\node[boxS,fill=whiteA] (SC) at (18.3,10.8)
		{偏离程度\quad $\mathrm{score}_i=\dfrac{\lvert x_i-m_i\rvert}{\mathrm{MAD}_i}$}; % 高度由 \dfrac 分数撑开
		
		\draw[arr] (A1.east) -- (B1.west);   % 水平箭头：盒子右缘→左缘
		\draw[arr] (A2.east) -- (B2.west);
		\draw[arr] (B1.east) -- (E.160);     % E.160/E.200：异常点集框左缘偏上/偏下的锚点
		\draw[arr] (B2.east) -- (E.200);
		\coordinate (J) at (12.9,10.8);      % 三线汇聚点
		\draw[pln] (W1.east) -- (J);
		\draw[pln] (W2.east) -- (J);
		\draw[pln] (W3.east) -- (J);
		\draw[arr] (J) -- (SC.west);
		
		% ================= Step 2：分析数据特征 =================
		\draw[subp,fill=cyanB] (23.8,8.0) rectangle (33.5,14.6); % 周期性子面板：宽 9.7，高 6.6
		\node[anchor=north west] at (24.2,14.35) {\textbf{周期性}};
		\node[boxS,fill=whiteA] (P1) at (28.7,12.7) {STL分解\quad $x_t=T_t+S_t+R_t$};
		\node[boxS,fill=whiteA] (P2) at (28.7,10.8) {季节振幅\; $A_s=\max(S_t)-\min(S_t)$};
		\node[boxS,fill=whiteA] (P3) at (28.7,9.0)  {TTF、ACF检验};
		\draw[arr] (P1.south) -- (P2.north); % 纵向箭头：盒子下缘→上缘
		\draw[arr] (P2.south) -- (P3.north);
		
		\draw[subp,fill=cyanB] (23.8,0.5) rectangle (33.5,7.5);  % 下降趋势子面板：宽 9.7，高 7.0
		\node[anchor=north west] at (24.2,7.25) {\textbf{下降趋势}};
		\node[boxS,fill=whiteA] (T1) at (28.7,5.1) {四季平均斜率 $\longrightarrow$ 季节性变化};
		\node[boxS,fill=whiteA] (T2) at (28.7,2.3) {年均劣化斜率 $\longrightarrow$ 跨年加速老化};
		
		% ================= Step 3：维护的影响 =================
		\node[boxS,fill=cyanC]  (G)  at (14.6,5.6) {增益系数G};
		\node[boxS,fill=cyanC]  (D)  at (20.3,5.6) {损伤系数 $\alpha$};
		\node[boxS,fill=greenC] (Mm) at (14.6,3.3) {中维护};
		\node[boxS,fill=greenC] (Bm) at (20.3,3.3) {大维护};
		\node[boxS,fill=cyanC]  (Se) at (17.5,1.0) {季节};
		\draw[arr] (G.south)  -- (Mm.north);  % 竖直箭头
		\draw[arr] (G.south)  -- (Bm.north);  % 交叉箭头
		\draw[arr] (D.south)  -- (Mm.north);  % 交叉箭头
		\draw[arr] (D.south)  -- (Bm.north);  % 竖直箭头
		\draw[arr] (Se.north) -- (Mm.south);  % 季节→中维护
		\draw[arr] (Se.north) -- (Bm.south);  % 季节→大维护
		
		% ================= Step 4：透水率指标 =================
		% 大括号：line width=2pt 括号线粗；amplitude=8pt 括号中点凸出宽度；
		% (1.0,6.7)--(1.0,0.5) 括号竖直长度 6.2cm；mirror 使凸尖朝左
		\draw[line width=2pt,decorate,decoration={brace,amplitude=8pt,mirror}]
		(1.0,6.7) -- (1.0,0.5);
		\node[anchor=west] at (1.7,5.9) % anchor=west：坐标为该行左端
		{自然衰减\quad $k_0=\dfrac{1}{M}\displaystyle\sum_{m=1}^{M}\lvert k_m\rvert$};
		\node[anchor=west] at (1.7,4.35)
		{季节影响\quad $\beta_s=\dfrac{\overline{k}_s}{k_0}$，\quad $s\in\{\text{春},\text{夏},\text{秋},\text{冬}\}$};
		\node[anchor=west,font=\normalsize] at (1.7,2.7) % 此行单独用 \normalsize 以容纳长公式
		{维护增益\quad $\displaystyle G=\frac{1}{n_a}\sum_{i=1}^{n_a}x\left(t_a^{(i)}\right)-\frac{1}{n_b}\sum_{i=1}^{n_b}x\left(t_b^{(i)}\right)$};
		\node[anchor=west] at (1.7,1.0)
		{维护损伤\quad $\alpha=\dfrac{k'_{i}-k'_{i-1}}{k'_{i-1}}$};
		
		% ================= 面板间流转箭头（坐标为箭头起、止点） =================
		\draw[arr] (22.85,12.3) -- (23.35,12.3);  % Step1 → Step2，长 0.5cm
		\draw[arr] (23.35,4.4)  -- (22.85,4.4);   % Step2 → Step3，长 0.5cm
		\draw[arr] (12.05,4.4)  -- (11.55,4.4);   % Step3 → Step4，长 0.5cm
		
	\end{tikzpicture}
\end{document}